\section{The Formalization and Its Key Decisions}
\label{sec:formalization}

The development is four Lean libraries.  \texttt{OracleComputability} fixes
a machine model of oracle computation and proves the facts any priority
argument over that model needs.  \texttt{FriedbergMuchnik} and
\texttt{SacksSplitting} are built on it, and neither imports the other;
Lake enforces this, so building either from a clean tree compiles no module
of the other.  \texttt{PriorityArguments} is a single root module joining
them.  The two headline theorems are:

\begin{lstlisting}
theorem friedberg_muchnik :
    exists A B : Set Nat,
      CE A /\ CE B /\ not (A <=T B) /\ not (B <=T A)

theorem sacks_splitting :
    forall A : Set Nat, CE A -> not (ComputableSet A) ->
      exists A0 A1, CE A0 /\ CE A1 /\ Disjoint A0 A1
        /\ A0 union A1 = A
        /\ not (A <=T A0) /\ not (A <=T A1)
\end{lstlisting}

The vocabulary --- \texttt{CE}, \texttt{ComputableSet}, and the relation
displayed as $\leT$ --- is defined in the shared library and imported by
both clients, not redefined, so the two theorems are stated in one
language.

\subsection{Deciding What Is Primary}

\paragraph{A finitary oracle model.}
The development does not start from a semantic type of oracle functionals.
Doing so would require an oracle-machine model, continuity or compactness
theorems, extraction of finite uses from semantic computations, and a
connection between functionals and executable indices --- all before the
first requirement could be stated.  Instead \emph{finite execution is
primary}.  The evaluator runs a concrete program for bounded fuel against a
partial oracle $\N \to \mathtt{Option}\ \mathtt{Bool}$, and returns one of
three outcomes: it halted, recording both its output and the use; it got
\emph{stuck}, having asked a question the finite oracle could not answer;
or it ran out of fuel.  Distinguishing the last two matters, because they
are repaired by different things --- more oracle information, or more
resources --- and a priority construction reacts to them differently.
Infinite computation is then defined as existence of a finite successful
run against a total oracle.

Recording the use inside the result, rather than recovering it afterwards
by a continuity argument, is the decision the rest of the development rests
on.  It makes the use principle a statement about data the construction can
inspect at a stage.  The principle itself, \texttt{run\_halt\_mono}, says
that a halting bounded run survives more fuel and any oracle that preserves
the answers it actually used.  The hypothesis is one-sided and deliberately
weak: positions at or above the recorded use are unconstrained, and so are
positions below the use where the original oracle had no information.  That
weakness is exactly what a restraint can deliver.

\paragraph{Finite sets, not Boolean prefixes.}
Evolving \ce.\ approximations are monotone finite sets, with Boolean lists
appearing only as bounded snapshots supplied to one computation:

\begin{lstlisting}
def StageMono (F : Nat -> Finset Nat) : Prop :=
  forall s, F s subset F (s + 1)

def limitSet (F : Nat -> Finset Nat) : Set Nat :=
  {n | exists s, n in F s}
\end{lstlisting}

The alternative --- representing an approximation as a prefix-monotone
Boolean string, where each bit is fixed once it appears --- is not merely
inconvenient, it proves the wrong theorem.  Fixing bits makes the limit
sets decidable, and the development proves that a decidable set reduces to
every oracle.  A prefix-monotone construction would therefore have handed
back precisely the reduction the theorem denies.  The two representations
coexist because they carry different information: a stage set may
accumulate new positive information at any later stage, while a snapshot
merely answers finitely many questions for one bounded run.

\paragraph{Where Mathlib enters, and where it does not.}
Ordinary, non-oracle computability is imported; oracle computability is
built locally.  Exactly three files import Mathlib's computability library:
one for the primitive-recursive closure theorems used to show the stage
functions computable, one embedding Mathlib's partial-recursive codes into
the query-free fragment of our language, and one relating our reducibility
to Mathlib's own oracle model.  The local oracle semantics --- the program
syntax, the evaluator, the use principle, the reducibility vocabulary, the
approximation layer --- mention Mathlib's codes nowhere, so changes to
Mathlib's internal coding decisions cannot reach them; the clients meet
Mathlib only where they must, in discharging computability of their stage
functions.  This
boundary lets the construction reason about explicit finite uses while
still reusing Mathlib's closure library \cite{carneiro2019,mathlibPartrec}.

\paragraph{Reducibility vocabulary.}
The central definition is that $A \leT B$ holds when some local oracle
program computes the characteristic function of $A$, on every input, from
the total characteristic function of $B$:

\begin{lstlisting}
def TuringReducible (A B : Set Nat) : Prop :=
  exists c : OracleCode,
    forall x : Nat,
      Computes c (charFun B) x (natOfBool (charFun A x))
\end{lstlisting}

Reflexivity, transitivity (by substituting one program for every query node
of another), and the fact that decidable sets reduce to every oracle are
all proved, so the relation is a preorder with the expected bottom
behaviour.  The program numbering is proved to be a bijection between $\N$
and the code type, which is what turns ``for every index $e$'' into ``for
every reduction'': the final diagonalization takes an arbitrary candidate
program, passes to its index, and lets the corresponding requirement kill
it.

\subsection{What Transferred to the Second Theorem, and What Did Not}
\label{sec:transfer}

The foundation was built for one theorem.  To test whether it was a
foundation rather than scaffolding, we formalized a second finite-injury
theorem against it and recorded, component by component and before writing
any of the second construction, what we expected to reuse.

\paragraph{The foundation transferred whole.}
Every component of the machine model, the use principle, the reducibility
vocabulary, the Mathlib bridge and the approximation layer was imported and
used unchanged.  Two items are worth naming.

The restraint lemma --- a computation seen halting against a stage snapshot
survives to the limit provided nothing below its recorded use is enumerated
later --- was already stated generically in the stage sequence rather than
in terms of the first construction's sets.  The preservation argument at
the heart of Sacks splitting therefore reuses it verbatim.  This was the
open question the reuse inventory was written to answer, and the answer is
clean.

More interestingly, the determinism lemma \texttt{run\_halt\_unique} carries
more weight in the second theorem than in the first.  In
Friedberg--Muchnik it says only that a computation has one output.  In
Sacks it is what makes \emph{the use of the limit computation at $y$} a
well-defined number, and that number is what bounds the restraint function.
Recording the use rather than only the output was a decision taken for the
first theorem, for reasons internal to it; it paid for a theorem it was not
designed for.  That is the strongest evidence we have that the foundation
generalizes.

\paragraph{The priority layer transferred not at all.}
The Friedberg--Muchnik injury argument counts \emph{events}.  A
requirement's record carries a rank in $\{0,1,2\}$ --- has it a witness,
has it acted --- receiving attention raises the rank strictly, and only
injury from above resets it.  So once the stronger requirements go quiet, a
requirement acts at most twice more.  This is the explicit progress measure
the informal proof does not supply, and it is small enough that the
induction along priority is arithmetic.

A splitting requirement generates no events.  It never acts: the elements
are not ours to choose, they arrive from the enumeration of the given set,
and the construction only decides which half each one joins.  Its restraint
moves continuously with the length of agreement between the reduction and
the set being split, so there is no discrete step to count and no rank to
raise.  The replacement is a three-part simultaneous induction along the
priority order: bounded restraints below $j$ give finite injury to $j$;
finite injury to $j$ gives satisfaction of $j$; satisfaction of $j$ bounds
$j$'s restraint.

The two arguments share a shape --- induction along the priority order, no
closed-form injury bound --- and no content.  The two files have no
declaration in common.  We looked for a common generalization worth stating
and did not find one, and we record that as a negative result about reuse
rather than as a defect of the first development's design: ``finite
injury'' names a family of arguments, not a single lemma one can prove once
and instantiate.

One further difference is forced by the mathematics rather than by taste.
The splitting restraint must be \emph{cumulative}, kept in the state as a
running maximum.  The raw stage restraint is not monotone --- the length of
agreement drops whenever a number below it enters the set --- and a
restraint that can drop protects nothing, since a computation frozen at one
stage could be destroyed later merely because its requirement had
temporarily stopped asking.  Friedberg--Muchnik needs no such device, its
restraint being written once, when a requirement acts.

\paragraph{Where the hypothesis is spent.}
Sacks splitting takes a hypothesis Friedberg--Muchnik does not have: the
set being split is not decidable.  The step that consumes it is the one
place where the second development required an argument with no analogue in
the first.  If a requirement is never injured after some stage and its
reduction nevertheless computes the characteristic function of $A$
everywhere, then $A$ is decidable: to decide $y \in A$, search for a stage
past the last injury whose length of agreement exceeds $y$, and report
whether $y$ is in the approximation at that stage.  Such a stage exists
because the reduction is total; the answer is correct because past the last
injury the agreeing computation is frozen by the restraint lemma, so it
\emph{is} the true value; and the search is effective because the whole
construction is primitive recursive.  Since $A$ is not decidable, no
requirement can be in that situation, so every requirement is satisfied.

Friedberg--Muchnik has no counterpart.  Its requirements are satisfiable
because their witnesses are fresh, which is a property of the construction;
a splitting requirement is satisfiable only because of a property of the
\emph{given} set.  We describe this as a step of the classical proof made
explicit rather than as new mathematics: it is the point at which a
textbook appeal to non-computability becomes a decision procedure that has
to be exhibited and shown effective.

\paragraph{What the second theorem cost.}
The shared foundation is roughly 2{,}870 lines; the Friedberg--Muchnik
client is roughly 2{,}000 and the Sacks client roughly 1{,}550.  The ratio
is the point: the expensive part --- model, evaluator, use principle,
Mathlib bridge, approximation layer --- was already present, and the second
theorem paid only for its own strategy.

Two components of the first development turned out to be artifacts of its
own design rather than necessities.  Its ten-field state invariant exists
to keep freely chosen witnesses fresh, distinct, and clear of
higher-priority restraints; a construction that does not choose its
elements needs none of it, and the two invariants that remain --- the
halves stay disjoint, and their union is the current stage of the given set
--- are one induction.  And its computability layer needs a shadow
implementation on plain tuples, because its state types are custom
structures with no \texttt{Primcodable} instance; defining the splitting
state as a tuple from the outset removes that layer, its simulation
theorem, and its list-surgery lemmas entirely.

The exercise also found two defects that only a second client could find.
Three generic helpers about primitive recursion had been declared
\texttt{private} inside the first development's computability file.
Nothing in them mentions either construction --- they are what any stage
function needs in order to be shown computable --- but \texttt{private}
made them invisible downstream.  Nothing needed generalizing; they needed
exporting.  Separately, the first development's root module imported a file
that did not exist, so a clean checkout did not build at all.  Its own
build never noticed, because the file was present locally.  Both are
invisible until someone tries to build on the code.

\subsection{Analytic, and Bridged Rather Than Asserted}
\label{sec:related}

Post's problem has been mechanized before, and the contrast is a design
decision rather than a survey obligation.  Zeng, Forster and Kirst
formalize a solution in Coq \cite{zeng2024post}, building on the synthetic
oracle computability of Forster, Kirst and M\"uck
\cite{forster2023oracle}.  Their setting is \emph{synthetic}: one works in
the Calculus of Inductive Constructions and takes every function
$\N\to\N$ to be computable by axiom, so there is no object-level machine
model and none of the details it brings.  That is a deliberate and
effective choice, and it answers the standard complaint that explicit
models of computation drown a computability proof in bookkeeping that stays
invisible on paper.  Two points of their work bear directly on ours: the
solution they formalize is Soare's low simple set, which their abstract
explicitly contrasts with ``the full Friedberg--Muchnik theorem
constructing two incomparable predicates''; and the same abstract records
that Bauer posed a synthetic proof of Friedberg--Muchnik as a challenge.

The present development is \emph{analytic}, and pays the cost the synthetic
approach avoids.  What is bought is that every informal appeal to
``preserve the computation below its use'' becomes an application of one
lemma about one evaluator, and that the resulting reducibility relation can
be compared with an external one by proof rather than by fiat.  Neither
approach dominates; they answer different questions about the same theorem.

That comparison is the last design decision.  Mathlib contains a model of
oracle computability \cite{mathlibTuringDegree}: an inductive relation
\texttt{RecursiveIn} on partial functions, Turing reducibility and
equivalence, and degrees as the quotient with a partial-order instance.  It
does not contain a use principle, oracle-relative step-indexed evaluation,
oracle-relative computable enumerability, or the jump, and it is not
connected to Mathlib's computably enumerable predicates, which live in a
separate file developed for the unrelativized theory.  The two libraries
are complementary: Mathlib has the degree \emph{structure} we do not build,
and we supply the finite-computation machinery it lacks.

So we prove the connection instead of assuming it.  A bridge file
establishes that every reduction in Mathlib's model is realized by a
program of our language, by induction over \texttt{RecursiveIn} with one
program per constructor.  This is the direction that matters: both of our
conclusions are \emph{negations} of reducibility, so the risk to guard
against is that our relation admits too few reductions and thereby makes
the negations cheap.  It does not, and both theorems are consequently
restated over Mathlib's relation.  The converse inclusion is not needed for
these results and is not proved.

Two priority claims, stated with the caution they deserve.  We are not
aware of another formalization of Post's problem in an analytic setting
with an explicit oracle machine model, nor of another in Lean; the
synthetic precedent above is the closest.  We are also not aware of any
formalization of the Sacks Splitting Theorem in any proof assistant.  Both
claims rest on a literature and repository search rather than an exhaustive
one, and neither is load-bearing for the technical content of this paper.

\subsection{The Shape of the Implementation}

The details of the machinery are tabulated in
Appendix~\ref{app:map}; what follows is the shape.

The program language is Kleene's schemata --- zero, successor, the two
projections, pairing, composition, primitive recursion and unbounded search
--- plus one oracle primitive, \texttt{query}, which asks whether its input
is in the oracle set.  A code with no \texttt{query} node is an ordinary
partial recursive program; the bridge embeds every Mathlib partial-recursive
code into that fragment and proves the embedding semantics-preserving, which
is the direction the construction consumes.  The evaluator's fuel discipline
deliberately copies
Mathlib's \texttt{evaln}, which is what lets the embedding be proved by a
single structural induction rather than by a simulation with fuel
translation.  The whole evaluator is proved primitive recursive, in the
form in which the construction consumes it: numbered codes in, encoded
results out.

The Friedberg--Muchnik construction keeps two enumeration lists, one record
per requirement holding a witness, an acted flag and a restraint, and a
freshness counter dominating everything mentioned so far.  At each stage
the least requirement requiring attention receives it: with no witness it
is given the counter's value; with an unacted witness whose diagonalizing
computation has appeared within the stage's resources, it acts ---
enumerating the witness, recording the computation's use as its restraint,
and resetting every weaker record.  A hypothesis-free case analysis proves
that every stage is exactly one of idle, appointment, or action, so no
execution path can escape the verification.  Ten invariants, proved
together by induction over those three transitions, carry the properties
the satisfaction proof cites by name, of which the load-bearing one says
that an acted requirement's diagonalizing computation still halts against
the current stage.

The splitting construction is smaller for a structural reason: it makes no
positive moves.  A stage takes the numbers the given set's own enumeration
has just produced and routes each to the half not being used as an oracle
by the strongest requirement that restrains it, defaulting to the first
half when nothing restrains it.  Two invariants suffice: the halves stay
disjoint, and their union is the current stage of the given set.

The computable-enumerability obligation is discharged the same way in both
developments and is the most expensive part of each.  The stage function is
shown primitive recursive using Mathlib's closure library, an unbounded
search over stages turns it into a partial recursive semi-decider, and the
bridge converts that into the project's \texttt{CE}.  For
Friedberg--Muchnik this requires the shadow tuple implementation and a
simulation theorem tying it to the real machine, and it dominates the build:
one file accounts for roughly seven eighths of the project's compile time.
The instructive failure along the way was assembling that derivation as a
single large combinator term, which made failed unifications fall back into
the pairing encodings and left the elaborator symbolically evaluating
arithmetic for tens of millions of heartbeats.  Confining each
tuple-sized comparison to its own lemma fixed it.
